\section{Conclusion} \label{sec:conclusion}

\subsection{Scientific Question Addressed} \label{sec:concl_question}
This work investigated whether exoplanet vetting pipelines can be rigorously evaluated and certified using a multi-dimensional validation protocol, rather than relying on standard single-split metrics. In doing so, we addressed a critical methodological gap: standard exoplanet vetting benchmarks typically optimize discrimination (AUROC) on homogeneous splits, which conflates physical likelihood (aleatoric uncertainty) with recovery stability (epistemic uncertainty), masking domain-shift collapse, calibration decay, and statistical redundancies. We formulated and demonstrated a systematic six-layer validation protocol (encompassing out-of-domain transfer, nested controls, cohort partitioning, noise stress-testing, and probability calibration) to evaluate exoplanet vetting systems.

\subsection{Supporting Evidence} \label{sec:concl_evidence}
We implemented and tested this validation protocol using the Transit Ambiguity Recovery System (TARS) and its graph-theoretic Recovery Ambiguity Index (RAI) as a baseline case study. The application of our protocol successfully diagnosed critical operational limits:
\begin{itemize}
\item \textbf{Out-of-Domain Transfer}: Exposed a severe performance collapse ($\text{AUROC} \approx 0.48$) and calibration decay (Kepler ECE $= 0.5212$) under space-mission domain shift, tracing the physical sensitivity limit to cadence-induced transit-edge blurring in Kepler long-cadence observations.
\item \textbf{Nested Control Redundancy}: Revealed via nested logistic regression controls (EXP-019) that the RAI lacks statistically significant independent predictive value once controlling for BLS SDE ($p = 0.090$), successfully detecting a feature redundancy that standard AUROC testing missed.
\item \textbf{Cohort Boundaries}: Isolated stellar convective noise on giant hosts ($\log g < 4.0$) as a physical boundary that deforms alias-graph topology, defining the operational envelope of the graph-theoretic classifier.
\item \textbf{Robustness Sweeps}: Verified that within its defined operating regime, the index is robust to high-frequency noise and missing cadences (up to 30\% data loss) while remaining invariant to systematic measurement bias.
\end{itemize}

\subsection{Core Scientific Conclusion} \label{sec:concl_core}
Within the evaluated space-mission datasets, we conclude that standard single-metric AUROC evaluations underrepresent vetting pipeline vulnerabilities. The multi-dimensional validation protocol proposed and tested in this work provides a reproducible, open-science framework for certifying exoplanet vetting pipelines. This work demonstrates that rigorous evaluation can systematically reveal hidden limitations and redundancies in exoplanet vetting algorithms that conventional single-metric benchmarks fail to detect, ensuring their scientific reliability before large-scale survey deployment.

\subsection{Future Research Program} \label{sec:concl_future}
This validation protocol establishes a foundation for certifying next-generation exoplanet search engines:
\begin{itemize}
\item \textbf{Framework Generalization}: Future work will apply this validation protocol to evaluate high-capacity deep learning models (such as 1D CNNs) and other independent search pipelines across Kepler, K2, and simulated PLATO datasets to map their cross-mission operational boundaries.
\item \textbf{Probabilistic Vetting Calibration}: We will investigate methods to dynamically recalibrate vetting probabilities under domain shifts, optimizing Radial Velocity (RV) and spectroscopic follow-up triage schedules.
\item \textbf{Graph-theoretic Enhancements}: We plan to extend the graph-theoretic baseline to incorporate Transit Timing Variations (TTVs) in multi-planet systems, adapt features for evolved host stars, and evaluate directed graphs to capture asymmetric timing uncertainties.
\end{itemize}
